\documentclass[11pt,a4paper]{article}
\usepackage{../shared/hanzocover}
\usepackage[utf8]{inputenc}
\usepackage[T1]{fontenc}
% Shared preamble for the compute-market series.
% One style, one vocabulary, inputted by every paper so the series reads as one voice.
\usepackage[margin=1in]{geometry}
\usepackage{booktabs}
\usepackage{amsmath}
\usepackage{amssymb}
\usepackage{amsthm}
\usepackage{graphicx}
\usepackage{xcolor}
\usepackage{hyperref}
\hypersetup{colorlinks=true, linkcolor=blue, urlcolor=blue, citecolor=blue}
\usepackage{enumitem}
\usepackage{listings}
\usepackage{array}
\usepackage{tabularx}
\usepackage{siunitx}
\sisetup{group-separator={,}, group-minimum-digits=4}

\definecolor{kw}{rgb}{0.13,0.29,0.53}
\definecolor{cmt}{rgb}{0.40,0.40,0.40}
\definecolor{str}{rgb}{0.16,0.50,0.16}
\lstset{
  basicstyle=\ttfamily\footnotesize,
  keywordstyle=\color{kw}\bfseries,
  commentstyle=\color{cmt}\itshape,
  stringstyle=\color{str},
  showstringspaces=false,
  breaklines=true,
  columns=fullflexible,
  frame=single,
  framesep=4pt,
}

\newtheorem{definition}{Definition}
\newtheorem{proposition}{Proposition}
\newtheorem{remark}{Remark}

\newcommand{\code}[1]{\texttt{#1}}
\newcommand{\repo}[1]{\texttt{#1}}

% Series vocabulary — used identically in every paper.
\newcommand{\job}{\ensuremath{J}}
\newcommand{\bundle}{\ensuremath{R}}
\newcommand{\cost}{\ensuremath{\mathcal{C}}}

\tolerance=800
\emergencystretch=3em


\title{Clearing Heterogeneous Compute}

\author{
Zach Kelling\thanks{Corresponding author: research@hanzo.ai} \\
\textit{Hanzo Industries Inc (Techstars '17)} \\
Los Angeles, California \\
\texttt{https://hanzo.ai}
}

\date{2026}

\begin{document}
\hanzocoverpage

\begin{abstract}
An exchange for compute is not an exchange for a commodity. The thing being
traded is a bundle --- device, memory, bandwidth, residency, locality, and the
right to a deadline --- and the bundles are not substitutes for each other except
under constraints that vary per job. Maintaining a separate order book per device
class throws away the structure that matters, and quoting a single price per
accelerator-hour throws away all of it.

This paper is the mechanics. We define the job and bundle types, state the
feasibility constraint, give the clearing objective including the residency
terms, and show why the natural objective of minimising price is wrong in a way
that is not marginal. We then define four order types --- including a rent on
state, which has no analogue in a stateless market --- describe how the KV cache
behaves as perishable inventory, and map each failure mode to a mechanism. We
close by identifying which parts of this we already run in open source and can
be lifted directly.
\end{abstract}

\tableofcontents
\newpage

\section{Jobs and bundles}

\begin{definition}[Job]
$\job = (m, c, \Delta, \kappa, \tau, V)$ with $m$ the model or computation
required, $c$ the context or input size, $\Delta$ the deadline, $\kappa$ the
determinism class, $\tau$ the required verification tier, and $V$ the buyer's
value at risk.
\end{definition}

\begin{definition}[Bundle]
$\bundle = (d, M, B, \beta, \alpha, \ell_{0})$ with $d$ the device, $M$ available
accelerator memory, $B$ ingress bandwidth, $\beta \in \{0,1\}$ whether the base
model for $m$ is resident, $\alpha \in \{0,1\}$ whether the buyer's adapter is
resident, and $\ell_{0}$ the network round-trip to the buyer.
\end{definition}

Feasibility is a hard constraint and must be evaluated before price:
\begin{equation}
\bundle \models \mathrm{cap}(\job)
\iff
M \ge |w_{m}| + |\mathrm{kv}(c)|
\;\wedge\;
c \le c_{\max}(d)
\;\wedge\;
\kappa \in \mathcal{K}(d)
\;\wedge\;
\tau \in \mathcal{T}(\bundle)
\;\wedge\;
\hat{t}(\job,\bundle) \le \Delta .
\end{equation}

The last conjunct is where residency enters feasibility rather than cost:
\begin{equation}
\hat{t}(\job,\bundle) \;=\; \ell_{0}
\;+\; (1-\beta)\frac{|w_{m}|}{B}
\;+\; (1-\alpha)\frac{|a|}{B}
\;+\; t_{\text{compute}} .
\end{equation}

\section{The objective, and why cheapest is wrong}

\begin{equation}
\bundle^{*} \;=\; \arg\min_{\bundle \,\models\, \mathrm{cap}(\job)}
\Big[
p(\bundle)\,t_{\text{compute}}
+ \lambda_{L}\ell
+ \lambda_{F}\phi
+ \lambda_{V}v(\tau)
+ \lambda_{D}\delta
+ \mu_{1}(1-\beta)
+ \mu_{2}(1-\alpha)
\Big],
\label{eq:clearing}
\end{equation}
with $\mu_{1}/\mu_{2} \approx 10^{3}$, the ratio of base-model size to adapter
size.

\begin{proposition}
Minimising $p(\bundle)$ alone is not an approximation to~\eqref{eq:clearing}; it
can select an infeasible bundle.
\end{proposition}

Take a job requiring a \SI{140}{\giga\byte} base under a \SI{3}{\second}
deadline. A device at \SI{0.50}{\per\hour} without the model resident, on a
\SI{10}{\giga\bit\per\second} link, needs \SI{112}{\second} to fetch the weights.
It is not expensive relative to a device at \SI{4}{\per\hour} holding the model
--- it is outside the feasible set, by a factor of \num{37} on the binding
constraint. Any allocator that ranks by price and then filters will spend most of
its time rejecting its own top choices.

Two practical consequences. Feasibility must be evaluated first, and residency
must be part of it. And $\mu_{1}$ is not a smooth penalty --- near an interactive
deadline it behaves as a step, which is why a market for interactive inference
looks less like a spot market and more like a market for scheduled capacity with
warm state.

\section{Order types}

\subsection{Intent}

The buyer states what it wants, not what to buy:
\begin{lstlisting}
intent:
  want:        "best answer to <request>"
  deadline:    4s
  max_price:   0.08
  min_tier:    T2
  determinism: D_inf
\end{lstlisting}
The exchange solves~\eqref{eq:clearing}. This is the order type agents should
use, and the reason is not convenience: an agent reasoning about device
specifications is an agent whose behaviour breaks when the fleet changes.

\subsection{Reserve}

A commitment to capacity over a window, priced below spot in exchange for the
provider's ability to plan. Standard, and it works here for the usual reasons.

\subsection{Rent on state}

New, and specific to this market. The buyer pays a continuous fee to keep its
adapter --- or its model, or a conversation prefix --- resident somewhere
capable.

\begin{equation}
\text{rent} \;=\; \rho \cdot |a| \cdot \mathrm{d}t
\qquad\text{buys}\qquad
\alpha = 1 \text{ with probability } \ge 1-\epsilon \text{ over the term.}
\end{equation}

The buyer converts a stochastic cold-start penalty into a predictable stream. The
provider receives payment for memory it would otherwise hold speculatively or
evict. This is an option on locality, it is continuously priceable, and it does
not exist in a stateless compute market because there is no state to rent.

For the provider it also resolves a real scheduling problem. An accelerator
holding $k$ resident adapters faces a knapsack decision on every arrival, and the
shadow price of a residency slot is the expected discounted value of its
occupant. Rent makes that shadow price observable instead of guessed.

\subsection{Search}

For Las Vegas workloads: a domain, a predicate, a success probability, a
deadline. Priced on expected work rather than measured time, settled on the first
valid witness. Treated fully in the companion paper.

\section{Cache as perishable inventory}

Prefix caching means a conversation's accumulated key-value state has value that
decays --- with time, as memory pressure evicts it, and with the conversation's
own probability of continuing. Writing $\gamma$ for the per-turn continuation
probability and $s$ for the saving from a cache hit, the expected value of
retaining a prefix for one more turn is $\gamma s$, against a holding cost equal
to the shadow price of the memory.

This is an inventory problem with a known solution shape, and it is the second
place --- after adapters --- where a continuous formulation earns its keep over
discrete order books. The quantity being priced varies smoothly, decays
predictably, and has a clean opportunity cost.

\section{Failure modes and mechanisms}

\begin{table}[h]
\centering
\begin{tabular}{lll}
\toprule
Failure & Mechanism & Note \\
\midrule
Provider defects mid-job & escrow, deadline refund & no adjudication needed \\
Deadline missed & automatic refund on timestamp & signing policy, not manual \\
Wrong output, class $\mathsf{D}$ & byte equality & proof of misbehaviour \\
Wrong output, class $\mathsf{S}$ & tier per job & buyer chooses \\
Residency misreported & self-revealing & the deadline is missed \\
Correlated committee & diversity as a constraint & otherwise $m$ collapses to 1 \\
Price manipulation & many providers, open quotes & thin markets stay thin \\
\bottomrule
\end{tabular}
\caption{Failure modes and the mechanism that addresses each.}
\end{table}

Residency deserves emphasis because it is unusually clean. A provider cannot
profitably claim to hold weights it does not hold, since the claim is falsified
by the deadline it then misses. The market needs no attestation for this, which
is rare and should be exploited rather than replaced with a certificate.

\section{What exists and can be lifted}

Two components are in production, open source, and directly reusable.

\paragraph{A Hamiltonian market maker.} The pricing layer of
equation~\eqref{eq:clearing} is implemented: roughly forty-seven hundred lines of
Rust carrying a hidden Markov core for regime detection, symplectic price
dynamics under a regime-dependent anharmonic potential, a compute-pricing adapter
that modulates by service-objective tightness and quality tier, per-resource
market states whose prices never cross, one-bit per-tenant preference adapters,
and expected-free-energy routing. Quotes are deterministic given a job and market
state. The companion paper on the market mechanism treats it in full; the point
here is that the objective above is not a proposal.

\paragraph{A learned per-request router over frontier arms.} The upper allocation
problem --- which intelligence, not which resources --- is solved by a router that
featurises a request in sub-microsecond time, scores every candidate arm by a
bilinear utility $u(x,p)=x^{\top}Wp$, and returns the argmax subject to hard
ceilings on latency, cost and quality, with the weight matrix fit offline by ridge
regression and adapted online per tenant by a contextual bandit. Its fan-out
variant selects the top-$N$ arms by the same utility, dispatches concurrently, and
fuses candidates with a synthesizer that has no weights of its own --- so a new
arm is onboarded by adding a profile row rather than retraining a policy over the
pool. Every request is metered on one plane, so cost is attributed per request to
the arm that served it. This runs in production behind a single endpoint.

The structural point for a market: this is the same shape as
equation~\eqref{eq:clearing} --- hard feasibility gating a soft-weighted objective
--- at a different layer and over a different feature space. Intelligence
allocation and capacity allocation are the same mechanism twice, and they compose
because their interfaces are narrow: the upper layer asks for capacity meeting a
specification without caring how price formed, the lower supplies a quote without
caring which model runs.

\paragraph{Open weights are what a provider network can actually serve.} A closed
frontier arm reachable only through a vendor API cannot be scheduled onto a
third-party accelerator; routing across such arms is a real product but the
provider network is not doing the work. Open-weight models are the substrate a
decentralized network can genuinely allocate --- serve anywhere with the memory,
fine-tune per tenant, distil, quantize. A market carrying both routes intelligence
at the top and allocates real compute underneath.

\paragraph{A model catalog with tiered routing.} Each model carries an ordered
list of route rungs. Each rung names an upstream, a provider, a context limit,
and two price pairs --- retail and cost --- so margin is explicit per rung rather
than implied. Alternate arms and a synthesis route are first-class. The catalog
validates its own invariants at load: a rung priced free may only route to a
zero-priced upstream, a rung may not be configured to serve below cost, and every
arm must name a known provider. Resolution maps a requested name to a concrete
model and identifier.

This is a working price-and-route layer over heterogeneous model supply. What it
does not yet price is state locality --- the $\mu$ terms of~\eqref{eq:clearing}
--- which is the extension this series argues for.

\paragraph{A request router.} Classification maps a request to a task from
extracted features --- text, approximate token count, media presence, explicit
task and modality hints --- kept pure so that a learned classifier can replace the
heuristic without touching policy. Downstream are policy, route, replica and
ring, giving pool selection, replica choice and consistent hashing.

The classifier boundary matters for an exchange: allocation quality depends on
correctly typing incoming demand, and a market that cannot improve its
classifier without redeploying its policy will not improve it.

\paragraph{An inference engine with adapters and a training API.} Low-rank and
quantized low-rank adapters, multi-adapter co-residency, mixture-of-adapters, and
a training loop exposed over the serving API. The gap --- per-sequence adapter
selection within a batch --- is stated precisely in the companion paper on
residency.

\paragraph{Settlement.} Threshold signing across twenty or more chains with
per-job escrow and no key reconstruction, behind a signer-agnostic assembler
boundary. Treated in the companion paper on custody.

\section{Sequencing}

The order that gets a market to real verified volume soonest, rather than the
order that looks most complete on a diagram: carry search volume first, because
verification is free and settlement is unambiguous. Add deterministic-class
inference next, because byte equality gives it verification at no marginal cost.
Add algebraic spot-checking for the general inference case. Add residency pricing
and state rent once adapters are per-request selectable, because until then the
$\mu$ terms are unobservable. Add succinct proofs last, for the statements that
settle on chain.

Every step in that order produces a market that works at the end of it, which is
the property worth optimising for.

\end{document}
